\begin{abstract}
Drifting models train a one-step generator by moving its samples along a field that attracts them to data and repels
them from each other. In text-to-image data, however, each caption usually comes with a single image. We show that
with one positive per condition the attraction term reduces exactly to $L_2$ regression onto that positive, so a
generator that cannot resolve nearby conditions collapses to the conditional mean. We propose \emph{Affinity Drifting}
(\affd), which replaces the single target by a kernel mean-shift over the positives of semantically similar conditions,
weighted by a caption-space affinity. For Gaussian kernels the field is the gradient of a log density ratio, which
lets us predict its fixed points: borrowing positives with self-only repulsion converges to a Nadaraya--Watson mixture,
while symmetric repulsion memorises. We test this bottom-up with pre-registered, single-factor experiments. In a
non-parametric particle system the predicted fixed points are reached at the resolution of the kernel. In a 2-D toy
whose conditional distributions are known, the original Drifting objective and the single-positive objective both
collapse to the conditional mean (mean-collapse index \EOneTwoDriftingMCI\ and \EOneTwoSingleMCI), while an oracle with
many positives does not (\EOneTwoOracleMCI). Borrowing positives restores the mixture: \affd covers every mode and
lowers the held-out $W_2$ to \EOneThreeRatio\ of the single-positive objective, matching the oracle. Controls show why:
averaging the borrowed targets without the $x$-dependent kernel still collapses, and random instead of semantic
neighbours spread samples off the conditional manifold.
\end{abstract}
